\chapter[Sermon 37. Whilst the vii. Seal Is Opened, and the Angels with Trumpets Come Forth, Christ the Intercessor of His Church Offers Up Before His Father the Prayers of His Faithful.]{Whilst the vii. Seal Is Opened, and the Angels with Trumpets Come Forth, Christ the Intercessor of His Church Offers Up Before His Father the Prayers of His Faithful.}
\label{sermon:037}
\markright{Sermon 37. Whilst the vii. Seal Is Opened, and the Angels with ...}

And when he had opened the seventh seal, there was silence in Heaven about the space of half an hour. And I saw seven Angels standing before God, and to them were given seven trumpets. And another Angel came and stood before the altar, having a golden censer, and much incense was given unto him, that he should offer the prayers of all saints upon the golden altar which was before the throne. And the smoke of the incense, which came from the prayers of all saints, ascended up before God from the hand of the Angel. And the Angel took the censer and filled it with fire from the altar, and cast it into the earth, and there were voices and thunderings and lightnings and an earthquake.

Surely there are no books in the world, written by whomever or whenever, which can compare with the books of holy Scripture, concerning sincere truth, pure simplicity, and plain order. And perhaps this is no marvel to anyone who knows that these writings were indeed written by men, but inspired by the Holy Spirit. There are buildings most skillfully constructed by men, framed and arranged in a most beautiful order. But what beauty will you judge them to have if you compare them with the creation of the world, and with that most beautiful order which we see daily in all created things, and the changeable course of times? The most excellent works of men have nothing in them, or seem vile, if you compare them with the workmanship of God the creator.

\index{Apocalypse (Book of)}But for the most brilliant order and most plain treatise, this book of the Apocalypse has among others a most notable, excellent, and wonderful praise. John promised some of the matter, signifying that he would speak of those things which should be done in the church from his time until the judgment. And the faithful know to what end they should take these things, not so that their curiosity might be maintained or satisfied, but that, being sufficiently warned beforehand, they should not fall, but take heed to themselves and hold fast to true salvation.

And inasmuch as there is much talk among men about why God does this or permits that, and why he does not prohibit these or those things, John has exhibited to us a most wholesome vision, by which we may learn not to talk against God, nor to contend with him, but to acknowledge all his judgments to be righteous and just. This thing, indeed, both all the saints in heaven and also angelic spirits acknowledge, and attribute all glory to God.

And thus, having prepared the minds of the listeners, he comes to the matter itself and declares the fatal destinies of the church. Under the sixth seal he touches generally the corruption of doctrine, which sin is more perilous and more pestilent than all dangers of the body or outward perils. He reasons more fully thereof, and now particularly under the opening of the seventh seal recites how far the same stretches. For he declares how many, how great, and what manner of sects, heresies, and troubles shall arise in the church, and how hurtful they shall be to the church.

\index{Judgment, Last}And this place contains a history of the corrupt doctrine, of heresies or sects, and troubles from the time of John until the last judgment; it is extended throughout the eighth, ninth, tenth, eleventh chapters.

Nevertheless, before the trumpets sound, for a consolation, as it were by a brief digression, a remedy is placed there, which the faithful in all ages may use in that pestiferous corruption to keep safe their souls and the integrity of the same. For many times in this book are brought in most strong consolations in matters of greatest difficulty. The tenth chapter will also serve this argument. And the remedy that he shows is this: that we must flee unto Christ, redeemer of mankind, intercessor and propitiator. And that we shall be safe under his defense, we must offer up to him our prayers continually. Verily, the Lord in the Gospel, reasoning of the greatest dangers of the Devil, and being at hand, yet adds by and by that which might comfort their sorrowful minds: “I have prayed for thee, Peter, that thy faith should not fail, \&c.” Behold, we are saved in greatest distress through Christ’s protection, that we should not faint in faith. However, as the Evangelical and Apostolic letters indicate everywhere, our continual prayers, which we offer to God through Christ, must be joined to our trust in Christ. And in few words, the intercession of Christ at the right hand of God, and the effect and manner of the prayer of the faithful are here set forth to behold.

But we shall declare everything in order. He spoke generally under the sixth seal of corrupt doctrine; in the seventh, he will declare the same particularly and most abundantly. And while the seventh seal was opened, there was silence in heaven almost half an hour. Concerning this silence, commentators write diversely. But as I think, this silence excites the hearers to a diligent and attentive hearing. For silence has an air of admiration and expectation of weighty matters. Solomon says in the ninth chapter of Ecclesiastes, “the words of wise men are heard in silence.” When weighty matters should be proclaimed and set forth, the herald is wont to proclaim silence. Indeed, the matters that follow are of great importance, and unless we observe them with great attentiveness, we shall perish in sects and seductions. Those spiritual wickednesses are more dangerous than corporal perils.

And now, while they remain silent, eagerly awaiting what will come, the last seal is opened. Behold, seven angels with trumpets appear, of whom we shall speak later.

Now a remedy is presented for these great evils, as I mentioned. To make it more vivid and impress it more deeply upon our hearts, it is shown through a divine vision. Before the throne, and almost encompassing it, appears a golden altar. And an angel comes and stands at this altar; he holds a golden censer into which the saints place their offerings. He offers them before the throne, and the fragrance of the incense ascends from the angel’s hand before God.

\index{Hebrews, Epistle to the}We said in another place that the golden altar of incense was the Lord Christ himself, who is both altar and sacrifice and priest, as Paul witnesses to the Hebrews. He is called an angel, namely, the one of whom Isaiah makes mention in chapter 9, and also Malachi, saying, “Behold, I send my angel, who shall prepare the way before me, and suddenly the Lord shall come unto his temple, whom you seek for, and the angel of the covenant, whom you desire: behold, he comes,” says the Lord of hosts. The former angel, that is to say, messenger or ambassador, was John the Baptist, who prepared the way for the Lord. He, namely, the later angel, came immediately after the preaching of John and made complete that everlasting covenant. He now appears on the right hand of God in Heaven.

And two things are said about him. First, that he stood before, or in, or upon the altar. We must not imagine anything physical here, but understand that this manner of speaking signifies the priesthood of Christ. He always appears before his Father on our behalf, as Paul taught in Romans 8 and Hebrews 9. Therefore, he pleads the cause of his church before God and is advocate for the faithful. The same one who stands before the altar also stands in the midst of the throne. For he is coequal with the Father in deity, according to which he sits on the throne; and in his humanity, he is of the same substance as we are, according to this dispensation, as both Bishop and true man, to stand before the altar. The latter, which is to be noted, is that Christ holds a golden censer in his hand. For he has taken our very nature without sin, so that he might intercede for us and offer our prayers to God the Father.

And lest any man should doubt that he receives our prayers and offers them to God, finally that the true office of the Church might also appear, offering up all things by Christ, there is added, “to him are given many fragrances.” But to what end? That he might give them upon the golden altar, and that before the seat, as though to say, that he might bring them into the sight of God.

And because of a further declaration, lest we should not know the true fragrances, which please God, and which the faithful offer unto God through Christ, one or twice he adds that those fragrances are the prayers of Saints. And he means by Saints, not those that dwell in heaven, but us in the earth, which are sanctified with the spirit of our God, with the blood of Christ, baptism, faith, and word. John 13. And the prayers are invocations and expressions of thanks. And he says expressly of all Saints, lest any should fear that he and his prayers offered by Christ were excluded: If thou believe, thou art holy, and thy prayer is of God accepted. What the prayers of Saints are, it appears in the Lord’s Prayer, which we offer up to the Father in the name and words of Christ: hallowed be thy name, thy kingdom come, and the rest, which all contend against those sects and corruptions of true doctrine.

\index{Irenaeus}Irenaeus refers to this passage in chapters 31 and 32 of the fourth book. By this, he calls the Eucharist—the giving of thanks—the sacrifice of Christians. Those who uphold papistry corrupt this passage, twisting it to mean that the priest should sacrifice the actual body of Christ for the living and the dead. But the holy Bishop of Lyons knew this foul error. Away with them and their sophistry, to where they deserve to be. I have also spoken before, somewhat about the same matter.

\index{Primasius}And that it might clearly appear unto all men, that the prayers of the faithful, offered to God through Christ, are pleasant and acceptable, there is added: and the smoke of the odours ascends, that is to say, the prayers of the faithful were accepted by God. Therefore let us offer diligently our prayers unto God through Christ. For he hears us, and delivers us from evil. And the scripture many times calls our prayers an acceptable sacrifice to God. The passages are in Hosea, 14, in the 50th Psalm. And in many other places. In Psalm 141, the prophet says, “Let my prayer be directed as incense in thy sight, the lifting up of my hands an evening sacrifice.” Primasius, expounding this place, said how Christ is said to have taken of the prayers of the saints. For because through him the prayers of all may come sweetly unto God. Hereof the Apostle: by him we offer up always a sacrifice of praise unto God, that is to say, the fruit of lips confessing his name.

Hereby is refuted the opinion of those who suppose that the saints in heaven are intercessors of the faithful, who should commend their prayers to God and make the way open to God. For what need have they to procure for themselves other intercessors or advocates? What lack do they find in Christ? Or whom may they prefer or compare with Christ? What shall we say that even now the offerings are offered by the hand of the angel? The celestial saints were present with the Lord and were seen about the throne, but which of them taking the censer and gathering the prayers of the faithful, offered them to God? It displeased King Azariah [? Ozias or Asarias] that he took in hand the censer, intending to sacrifice and to perform the priests’ office; the same would be worse for those who dwell on earth; indeed, they would not remain in heaven if they took upon themselves the office of the only Bishop. \&c.

After this, we have heard that Christ filled the censer with fire taken from the altar, and sent it down into the earth. By this narration, he returns to finish the exposition of the trumpets. This fire is the grace of the Holy Spirit, which is put into the censer, taken from the altar, and sent down into the earth. For Christ took the fullness of the Spirit, as John shows in the 1st and 3rd chapters. Christ is altar and censer. Of the altar, fire is taken. For the Holy Spirit is the spirit of the Father and of the Son. “Him,” he says, “I will send you from my Father.” He sent him into the earth under the shape of fiery tongues; he sends him also at this day into the hearts of the faithful, that he may inflame them. This is the same fire which the Lord in the Gospel of Luke says he will send into the earth, and would that it should burn.

Moreover, the effect of this fire follows immediately. For there were thunderings, and voices, and lightning, and earthquakes. By the voices of the Gospel, the wounds of sinners are healed, and the hearts of men enlightened by the illumination of the Holy Spirit, and so forth. Of these things we have also spoken in chapter 4, sermon 24. Concerning the preaching of the Gospel, as Haggai also prophesied it would come to pass, it sows a wonderful commotion of all nations, and so forth. Satan also was stirred, who raised up his ministers throughout the world against the wholesome preaching of the Gospel. For sects sprang up, whom the defenders of the truth resisted, fighting with them. Whereof now he will reason at length. The Lord grant grace that these things may both be spoken and heard with much fruit.
